\section{Method}
This semester 12 MSc students has been working with the Naiad project, as supposed to the 18 MSc students working with it the previous year. The new group was divided into smaller teams by discussing every members background and interests. The decision fell on five groups, three groups working with software and two with hardware. For the software groups, one was working with the vision system, another with the graphical user interface and the third with the Naiad main control software. The hardware group was divided into two smaller groups where one focused its work on electronics and the other on mechanics. The three software groups was managed by a software manager and the hardware groups by a hardware manager. The hierarchy is shown in fig. \ref{hierarchy}.
\begin{figure}[!ht]
	\begin{center}
		\includegraphics[width=120mm]{./Images/Method/hierarchy.png}
		\caption{The hierarchy of the project.}
		\label{hierarchy}
	\end{center}
\end{figure}

The project was managed in a scrum like method where different tasks were handed out using a web-based system called Trello\cite{trello}. Here the team members could easily access and see which tasks where available. The tasks where prioritized using a five grade scale, where a low number meant a high priority.
\begin{enumerate}
\item Crucial, needed for operating in water.
\item Important, needed for maintaining control in water. 
\item Necessary, needed for entering RoboSub.  
\item Good to have, needed to win RoboSub. 
\item Not required, needed for other missions than RoboSub. 
\end{enumerate}
All of the requests from the client was prioritized together with the whole group. Files and documents was shared using a cloud service called Dropbox\cite{dropbox}. 

The project was managed in a way where the team members were trusted to solve the problems in their own way. The project manager worked as a bridge between the groups trying to motivate the team members towards the same goal. The administrative work, as well as most of the sponsorship work, was  conducted by the project manager. The group managers did the administrative work of the groups, which included weekly reports to the project manager on the progress in the group as well as prioritizing the work.

Since the project was already started when this group took over, some time was put into getting familiar with the previous work. Some of the previous work had to be redone to better suit the purpose and current goals of Naiad. Almost all of the constructed hardware was kept in its previous state, however changes were made both to software and to some of the hardware designs that were yet to be implemented.